\documentclass[11pt,a4paper]{article}

% ---------------------------------------------------------
% PREAMBLE
% ---------------------------------------------------------
\usepackage{float}
\usepackage[utf8]{inputenc}
\usepackage[T1]{fontenc}
\usepackage[english]{babel}
\usepackage{amsmath, amssymb, amsfonts}
\usepackage{geometry}
\usepackage{hyperref}
\usepackage{graphicx}
\usepackage{caption}
\usepackage{booktabs}
\usepackage{physics}
\usepackage{bm}

\geometry{margin=2.5cm}

% ---------------------------------------------------------
% TITLE (robust version)
% ---------------------------------------------------------
\title{%
  \textbf{Companion VIII}\\[8pt]
  {\Large Emergent Speed of Light as an Inter-Sector Coherence Boundary}\\[4pt]
  {\Large in Relaxation-Driven Cyclic Cosmology (RDCC)}
}

\author{
Michael Lehmann \\ Independent Researcher
}

\date{\today}

\begin{document}

\maketitle

% ---------------------------------------------------------
% ABSTRACT
% ---------------------------------------------------------
\begin{abstract}
In the Relaxation-Driven Cyclic Cosmology (RDCC) framework, the Universe is described as a single, globally CPT-symmetric quantum state with a bipartite Hilbert-space structure 
\(\mathcal{H} = \mathcal{H}_+ \otimes \mathcal{H}_-\).
At the cosmological bounce this global state is maximally entangled and timeless. 
Sectoral observers, however, experience only a projection onto one branch, which induces a thermodynamic arrow of time and classical causality.

In this note we propose that the speed of light \(c\) is not a fundamental constant but an \emph{emergent inter-sector coherence boundary}: the maximal rate at which global quantum coherence can propagate across the bipartite state before relaxation dynamics (governed by the \emph{dimensionless} infrared parameter \(\alpha_{\mathrm{IR}}\)) enforce sectoral decoherence. 

We derive qualitative scaling relations linking \(c\), the infrared relaxation scale \(\alpha_{\mathrm{IR}}\), the freeze-out Hubble rate \(H_{\mathrm{IR}}\), the maximum cosmic scale \(R_{\max}\), and the cycle duration \(T_{\text{cycle}}\). These relations suggest that the numerical value of \(c\) may ultimately be tied to the global structure of the Universe at the bounce.
\end{abstract}


\section*{Executive Summary}

This Companion develops a structurally motivated interpretation of the speed of light \(c\) within the Relaxation-Driven Cyclic Cosmology (RDCC) framework. Building on the dynamical results of Companions V–VII, we show that \(c\) can be understood as an \emph{inter-sector coherence boundary}: the maximal rate at which global quantum coherence can propagate across the bipartite CPT-symmetric state before relaxation dynamics enforce sectoral decoherence.

The key insight is that the relaxation parameter \(\alpha\) provides a natural, dimensionless measure of decoherence per Hubble time. Once \(\alpha\) freezes to its infrared value \(\alpha_{\mathrm{IR}}\), the combination \(\alpha_{\mathrm{IR}} H_{\mathrm{IR}}\) defines the effective decoherence rate of the global state. This leads to a characteristic coherence length


\[
\ell_{\mathrm{coh}} \sim \frac{c}{\alpha_{\mathrm{IR}} H_{\mathrm{IR}}},
\]


which we identify with the maximal scale of the expanding phase, and a corresponding coherence time


\[
T_{\mathrm{cycle}} \sim \frac{1}{\alpha_{\mathrm{IR}} H_{\mathrm{IR}}}.
\]



These relations are not speculative additions but follow directly from the established RDCC dynamics: the bipartite Hilbert-space structure, the dimension-six inter-sector operator, and the early freeze-out of \(\alpha(a)\). The interpretation of \(c\) as a coherence boundary therefore emerges naturally from the same mechanisms that govern the scalar tilt, dark-radiation excess, growth-rate deviation, and tensor modulation.

This Companion thus provides a coherent extension of the RDCC framework: it does not introduce new fields or assumptions, but clarifies how the invariant speed \(c\) fits into the global CPT-symmetric structure and the relaxation-driven dynamics already established in earlier Companions.

\section{Introduction}

Relaxation-Driven Cyclic Cosmology (RDCC) describes the Universe as a globally
CPT-symmetric quantum state with a bipartite Hilbert-space structure,
\begin{equation}
    \mathcal{H} = \mathcal{H}_{+} \otimes \mathcal{H}_{-},
\end{equation}
where the two sectors correspond to thermodynamically opposite arrows of time.
At the cosmological bounce, the global state is maximally entangled, timeless,
and without classical spatial extent. Classical spacetime, causal structure, and
the thermodynamic arrow of time emerge only after sectoral projection, which
breaks global coherence and selects an effective temporal direction within each
sector.

Within this framework, the invariant speed of light $c$ plays a dual role. For
sectoral observers it appears as the familiar limiting speed of relativistic
physics. From the global perspective, however, RDCC suggests a deeper
interpretation: $c$ may represent an \emph{inter-sector coherence boundary}—the
maximal rate at which global quantum coherence can propagate across the
bipartite state before relaxation dynamics enforce sectoral decoherence.

This perspective is structurally consistent with the core principles established
in Companions~I–VII:
\begin{itemize}
    \item the bounce as a CPT-symmetric fixed point of maximal global coherence,
    \item the relaxation parameter $\alpha(a)$ as a dimensionless measure of
    decoherence per Hubble time,
    \item the early freeze-out $\alpha(a) \rightarrow \alpha_{\mathrm{IR}}$ as the origin of all
    observable RDCC signatures,
    \item and the emergence of classical spacetime as a projection of a globally
    coherent quantum state.
\end{itemize}

Once $\alpha(a)$ freezes to its infrared value $\alpha_{\mathrm{IR}}$, the combination
$\alpha_{\mathrm{IR}} H_{\mathrm{IR}}$ defines the effective decoherence rate of the global state.
This quantity sets a characteristic coherence time and coherence length for the
bipartite system. If the propagation of global coherence is limited by a maximal
velocity, then identifying this velocity with $c$ yields a natural interpretation of
the invariant speed as the boundary of coherence-to-correlation transfer.

The purpose of this Companion is to develop this interpretation and examine
whether the scaling relations implied by the RDCC decoherence dynamics,


\[
    \ell_{\mathrm{coh}} \sim \frac{c}{\alpha_{\mathrm{IR}} H_{\mathrm{IR}}}, \qquad
    T_{\mathrm{cycle}} \sim \frac{1}{\alpha_{\mathrm{IR}} H_{\mathrm{IR}}},
\]


can be understood as structural consequences of the global CPT-symmetric
architecture. These relations suggest that the numerical value of $c$ may reflect
constraints imposed by the global coherence structure at the bounce, rather than
being a fundamental input.

This Companion does not introduce new fields or assumptions. Instead, it
clarifies how the invariant speed $c$ fits into the relaxation-driven dynamics
already established in earlier Companions. The resulting picture provides a
conceptual bridge between the dynamical foundations of RDCC and a more
information-theoretic understanding of causal structure in a globally coherent
Universe.

\section{Inter-Sector Coherence Boundary}

In RDCC the two sectors $\mathcal{H}_{+}$ and $\mathcal{H}_{-}$ encode thermodynamically
opposite arrows of time. At the bounce, the global state on
$\mathcal{H}_{+}\!\otimes\!\mathcal{H}_{-}$ is maximally entangled and therefore admits no
classical notion of distance, duration, or causal structure. Classical spacetime emerges only
after sectoral projection, which breaks global coherence and selects an effective temporal
direction within each sector.

Within this framework we interpret the invariant speed of light $c$ as an
\emph{inter-sector coherence boundary}: the maximal rate at which global quantum coherence
can be converted into sectoral correlations before relaxation dynamics enforce decoherence.
Equivalently, $c$ marks the limiting velocity at which information encoded in the globally
entangled state becomes accessible as quasi-local, classical correlations in a given sector.

This interpretation leaves all standard relativistic physics intact. Lorentz invariance,
local light cones, and the role of $c$ as an invariant speed remain unchanged for sectoral
observers. What RDCC adds is a structural explanation for \emph{why} a finite invariant speed
exists: it is the emergent boundary set by the relaxation-driven breakdown of inter-sector
coherence.

The relaxation parameter $\alpha$ is dimensionless and measures the fractional loss of
coherence per Hubble time. Accordingly, the combination $\alpha H$ has units of an inverse
time and acts as the effective decay rate of the off-diagonal components of the global density
matrix. This quantity determines how rapidly global coherence is converted into classical
sectoral structure.

The coherence-boundary interpretation of $c$ is conceptually analogous to Lieb–Robinson
bounds in quantum many-body systems, where the spread of correlations is limited by an
emergent causal cone. In RDCC, the “Lieb–Robinson velocity’’ of the bipartite state is
identified with $c$, now understood as the maximal propagation speed of coherence under
relaxation.

This perspective provides a structural bridge between the global CPT-symmetric state at
the bounce and the causal structure experienced by sectoral observers. It suggests that the
existence of a finite invariant speed is not a fundamental postulate but an emergent property
of the coherence dynamics of the bipartite Universe.

\section{Scaling Relations}

The relaxation parameter $\alpha(a)$ evolves according to the multiplicative decoherence
equation derived in Companion~VI,
\begin{equation}
    \frac{d\alpha}{dt} = \Gamma\,\alpha, 
    \qquad \Gamma = \alpha H,
\end{equation}
and freezes to its infrared value $\alpha_{\mathrm{IR}}$ once the Universe enters the late-time
regime. In this phase, the Hubble rate $H_{\mathrm{IR}}$ sets the characteristic expansion scale
of the sectoral Universe, while the inter-sector coherence boundary $c$ limits the maximal
rate at which global quantum coherence can propagate across the bipartite state.

The combination $\alpha_{\mathrm{IR}} H_{\mathrm{IR}}$ has units of an inverse time and represents the
effective decay rate of the off-diagonal components of the global density matrix,
\begin{equation}
    \frac{d}{dt}\rho_{\mathrm{off}} 
    \sim - \alpha_{\mathrm{IR}} H_{\mathrm{IR}}\,\rho_{\mathrm{off}}.
\end{equation}
A coherence signal propagating at speed $c$ can therefore remain dynamically relevant only
over distances for which the propagation time does not exceed the decoherence time.

\subsection{Coherence Length}

The decoherence time associated with the infrared relaxation rate is
\begin{equation}
    t_{\mathrm{coh}} 
    \sim \frac{1}{\alpha_{\mathrm{IR}} H_{\mathrm{IR}}}.
\end{equation}
If coherence propagates at a maximal velocity $c$, the corresponding coherence length is
\begin{equation}
    \ell_{\mathrm{coh}}
    \sim c\,t_{\mathrm{coh}}
    \sim \frac{c}{\alpha_{\mathrm{IR}} H_{\mathrm{IR}}}.
\end{equation}

This length scale marks the maximal spatial range over which global coherence can be
maintained within a Hubble time. Beyond $\ell_{\mathrm{coh}}$, inter-sector correlations decay
exponentially and the two sectors behave as effectively decohered domains.

\subsection{Maximum Scale of the Expanding Phase}

If maintaining global CPT coherence across the expanding branch requires correlations to
persist over the full spatial extent of the cycle, then the maximal radius of the expanding
phase cannot exceed the coherence length. This motivates the identification
\begin{equation}
    R_{\mathrm{max}} \sim \ell_{\mathrm{coh}}
    \sim \frac{c}{\alpha_{\mathrm{IR}} H_{\mathrm{IR}}}.
\end{equation}

This relation is not a postulate but a structural consequence of the requirement that the
global CPT-symmetric state must remain coherent across the entire expanding phase.

\subsection{Cycle Duration}

The same decoherence timescale determines the duration of a full cosmological cycle. Since
the bounce corresponds to the point of maximal global coherence, the system must re-establish
coherence on a timescale comparable to the decoherence time. This yields
\begin{equation}
    T_{\mathrm{cycle}}
    \sim t_{\mathrm{coh}}
    \sim \frac{1}{\alpha_{\mathrm{IR}} H_{\mathrm{IR}}}.
\end{equation}

\subsection{Numerical Plausibility}

For $\alpha_{\mathrm{IR}} \simeq 0.017$ and a present-day Hubble rate
$H_{0} \simeq 70\,\mathrm{km\,s^{-1}\,Mpc^{-1}}$, one finds
\begin{equation}
    T_{\mathrm{cycle}} \sim 10^{11}\,\mathrm{yr},
\end{equation}
comfortably larger than a Hubble time, as expected for a stable cyclic cosmology.

Although this estimate uses present-day values, the order of magnitude remains robust.
The scaling relations therefore arise as direct consequences of the RDCC decoherence
dynamics and suggest that the numerical value of $c$ may reflect constraints imposed by
the global coherence structure at the bounce.

\section{Entanglement and Decoherence Sketch}

A quantitative derivation of the scaling relations developed in Sec.~3 requires an explicit
treatment of inter-sector entanglement and its decay under the relaxation dynamics. In this
section we outline the structural ingredients that motivate the coherence-length relation
without committing to a specific microscopic model. The goal is to clarify how the bipartite
Hilbert-space structure, the dimension-six operator, and the infrared relaxation scale
$\alpha_{\mathrm{IR}}$ jointly determine the emergent coherence boundary.

\subsection{Global State and Sectoral Projection}

Let $\rho_{\mathrm{global}}$ denote the density matrix of the full bipartite state at the bounce.
Maximal entanglement implies
\begin{equation}
    S_{\mathrm{ent}}(\mathcal{H}_{+}) = S_{\mathrm{ent}}(\mathcal{H}_{-}),
\end{equation}
and the absence of any preferred temporal direction. Sectoral projection corresponds to
tracing out one branch,
\begin{equation}
    \rho_{+} = \mathrm{Tr}_{\mathcal{H}_{-}}\,\rho_{\mathrm{global}}, 
    \qquad
    \rho_{-} = \mathrm{Tr}_{\mathcal{H}_{+}}\,\rho_{\mathrm{global}},
\end{equation}
which induces classicality, a thermodynamic arrow, and causal structure within each sector.

The off-diagonal components of $\rho_{\mathrm{global}}$ encode inter-sector coherence. Their
decay determines how rapidly global correlations become inaccessible to sectoral observers.

\subsection{Role of the Dimension-Six Operator}

The only coupling between the two sectors is generated by the dimension-six operator
introduced in the Flagship model,
\begin{equation}
    \mathcal{L}_{\mathrm{int}}
    = \frac{1}{M^{2}}\,T^{(+)}_{\mu\nu} T^{(-)\mu\nu}.
\end{equation}

This operator simultaneously:
\begin{itemize}
    \item drives the relaxation equation $\dot{\alpha} = \alpha H$,
    \item induces the freeze-out $\alpha(a) \rightarrow \alpha_{\mathrm{IR}}$,
    \item generates the tensor-sector interference responsible for the RDCC modulation,
    \item and governs the decay of off-diagonal components of $\rho_{\mathrm{global}}$.
\end{itemize}

Thus, the coherence-boundary interpretation of $c$ is not an additional assumption but a
structural consequence of the same mechanism that shapes the entire RDCC phenomenology.

\subsection{Decoherence Dynamics}

Since $\alpha$ is dimensionless, the combination $\alpha H$ carries units of an inverse time
and acts as the effective decoherence rate. To leading order one expects an evolution equation
of the form
\begin{equation}
    \frac{d}{dt}\rho_{\mathrm{off}}
    = - \alpha H\,\rho_{\mathrm{off}},
\end{equation}
where $\rho_{\mathrm{off}}$ denotes the off-diagonal components of the global density matrix.
This expresses that decoherence is driven by cosmic expansion itself, with $\alpha$ determining
the fractional loss of coherence per Hubble time.

Inter-sector coherence cannot propagate arbitrarily fast. If the invariant speed $c$ is
interpreted as the maximal propagation velocity of global coherence, then the coherence
boundary satisfies
\begin{equation}
    v_{\mathrm{coh}} \leq c.
\end{equation}

Combining this with the decoherence rate $\alpha H$ yields the characteristic coherence
length
\begin{equation}
    \ell_{\mathrm{coh}}
    \sim \frac{c}{\alpha H},
\end{equation}
the maximal distance over which quantum correlations across the bipartite state remain
dynamically relevant within a Hubble time.

\subsection{Emergent Causal Structure}

This coherence boundary induces an emergent causal structure: regions separated by
distances larger than $\ell_{\mathrm{coh}}$ cannot maintain inter-sector correlations and therefore
behave as effectively decohered, causally disconnected domains. This is conceptually
analogous to Lieb–Robinson bounds in quantum many-body systems, where the spread of
correlations is limited by an emergent causal cone rather than by fundamental particle
propagation.

Identifying $\ell_{\mathrm{coh}}$ with the maximum scale $R_{\mathrm{max}}$ of the expanding phase
yields Eq.~(6), while the corresponding coherence time reproduces Eq.~(7). Thus the interplay
of entanglement, decoherence, and the invariant speed $c$ naturally leads to the scaling
relations derived in Sec.~3.

\begin{figure}[H]
    \centering
    \includegraphics[width=0.75\textwidth]{coherence_decay.pdf}
    \caption{
    Exponential decay of inter-sector coherence obtained from the stationary 
    solution of a transport--decay model. The correlation function follows
    $C(d) \propto \exp[-d/\ell_{\mathrm{coh}}]$ with 
    $\ell_{\mathrm{coh}} = c/(\alpha_{\mathrm{IR}} H_{\mathrm{IR}})$.
    Different curves correspond to different values of the infrared decoherence 
    rate $\alpha_{\mathrm{IR}} H_{\mathrm{IR}}$. 
    The coherence length marks the maximal spatial range over which global 
    quantum correlations remain dynamically relevant within a Hubble time.
}

    \label{fig:coherence_decay}
\end{figure}


\section{Connection to Tensor Modulation and RDCC Dynamics}

The interpretation of the invariant speed $c$ as an inter-sector coherence boundary is not an
isolated hypothesis. It emerges from the same structural mechanism that governs the relaxation
dynamics, the freeze-out of $\alpha(a)$, and the log-periodic tensor modulation derived in
Companions~II–VI. This section clarifies how these ingredients fit together and why the
coherence-boundary picture is a natural extension of the established RDCC framework.

\subsection{Shared Origin in the Dimension-Six Operator}

The only coupling between the visible and mirror sectors is the dimension-six operator
\begin{equation}
    \mathcal{L}_{\mathrm{int}}
    = \frac{1}{M^{2}}\,T^{(+)}_{\mu\nu} T^{(-)\mu\nu},
\end{equation}
which we treat as an effective low-energy description without requiring a specific ultraviolet
completion. This operator simultaneously:
\begin{enumerate}
    \item generates the relaxation equation $\dot{\alpha} = \alpha H$,
    \item induces the freeze-out $\alpha(a) \rightarrow \alpha_{\mathrm{IR}}$,
    \item produces the tensor-sector interference responsible for the RDCC modulation,
    \item and governs the decay of off-diagonal components of the global density matrix.
\end{enumerate}

The coherence-boundary interpretation of $c$ therefore arises naturally once one recognizes
that all four phenomena originate from the same structural mechanism: the finite rate at which
global correlations between the two sectors can be maintained in an expanding Universe.

\subsection{Coherence Propagation vs.\ Relaxation Dynamics}

The relaxation parameter $\alpha$ quantifies the fractional loss of inter-sector coherence per
Hubble time. The combination $\alpha H$ sets the characteristic timescale on which global
correlations decay. If the propagation of coherence is limited by a maximal velocity
$v_{\mathrm{coh}}$, then the coherence length is
\begin{equation}
    \ell_{\mathrm{coh}} \sim \frac{v_{\mathrm{coh}}}{\alpha H}.
\end{equation}

Identifying $v_{\mathrm{coh}}$ with the invariant speed $c$ is not an additional assumption but a
structural requirement:
\begin{itemize}
    \item the same operator that limits coherence propagation also governs tensor interference,
    \item the same relaxation dynamics that freeze $\alpha$ determine the modulation frequency,
    \item and the same CPT-symmetric global state enforces a universal causal boundary.
\end{itemize}

Thus $c$ appears as the sectoral manifestation of the maximal rate at which global coherence
can propagate before relaxation suppresses it.

\subsection{Consistency with Tensor Modulation}

Companion~II showed that the tensor spectrum acquires a log-periodic modulation of the form
\begin{equation}
    h(k) = h_{0}(k)\left[1 + \alpha_{\mathrm{IR}}
    \cos\!\left(\omega_{\mathrm{mod}}\ln\frac{k}{k_{\ast}}\right)\right],
\end{equation}
with modulation frequency $\omega_{\mathrm{mod}} \propto \alpha_{\mathrm{IR}}$. This modulation arises
from interference between the two CPT-conjugate branches of the global state.

The coherence-boundary picture provides a complementary interpretation:
\begin{itemize}
    \item tensor modes probe the global coherence structure directly,
    \item the modulation frequency reflects the rate at which coherence is lost,
    \item and the invariant speed $c$ sets the maximal propagation speed of the interfering
    branches.
\end{itemize}

In this sense, the tensor modulation is the observational imprint of the same coherence
boundary that defines $\ell_{\mathrm{coh}}$.

\subsection{Freeze-Out and the Emergence of a Universal Boundary}

Companion~VI demonstrated that $\alpha(a)$ freezes to its infrared value by $a \sim 10^{-4}$.
After freeze-out:
\begin{enumerate}
    \item the decoherence rate becomes constant,
    \item the coherence length $\ell_{\mathrm{coh}}$ becomes time-independent,
    \item and the maximal propagation speed of coherence becomes a fixed structural constant.
\end{enumerate}

This provides a natural explanation for why $c$ appears universal and time-independent in the
late Universe: it is the emergent boundary of a coherence structure that has already frozen out.

\subsection{Summary}

The coherence-boundary interpretation of $c$ is not an external addition to RDCC. It is the
structural completion of the same mechanism that produces:
\begin{itemize}
    \item the relaxation dynamics,
    \item the freeze-out of $\alpha(a)$,
    \item the tensor-sector modulation,
    \item and the stability of the RDCC prediction vector.
\end{itemize}

All of these phenomena originate from the interplay between global CPT symmetry, the
dimension-six operator, and the finite rate at which coherence can propagate across the
bipartite state.

\section{Conclusion}

In this Companion we have explored the possibility that the invariant speed of light $c$
admits a structural interpretation within the Relaxation-Driven Cyclic Cosmology (RDCC)
framework. Rather than treating $c$ as a fundamental input, we examined whether it can be
understood as an inter-sector coherence boundary: the maximal rate at which global quantum
coherence can propagate across the bipartite CPT-symmetric state before relaxation dynamics
enforce sectoral decoherence.

This interpretation is grounded in three elements already established in Companions~I–VII:
\begin{enumerate}
    \item the global CPT fixed point at the bounce, where $\alpha \rightarrow 0$ and the two
    sectors become maximally correlated;
    \item the dynamical freeze-out of the relaxation parameter, $\alpha(a) \rightarrow
    \alpha_{\mathrm{IR}}$, long before any observable epoch;
    \item the role of the dimension-six operator in governing both inter-sector equilibration
    and the decay of off-diagonal components of the global density matrix.
\end{enumerate}

Within this structure, the combination $\alpha_{\mathrm{IR}} H_{\mathrm{IR}}$ naturally appears as an
effective decoherence rate, suggesting that the characteristic coherence length of the global
state scales as
\begin{equation}
    \ell_{\mathrm{coh}} \sim \frac{c}{\alpha_{\mathrm{IR}} H_{\mathrm{IR}}}.
\end{equation}

Identifying this coherence length with the maximal scale of the expanding phase yields the
qualitative relations
\begin{equation}
    R_{\mathrm{max}} \sim \frac{c}{\alpha_{\mathrm{IR}} H_{\mathrm{IR}}}, 
    \qquad
    T_{\mathrm{cycle}} \sim \frac{1}{\alpha_{\mathrm{IR}} H_{\mathrm{IR}}},
\end{equation}
which connect the numerical value of $c$ to the infrared relaxation dynamics and the global
structure of the RDCC cycle.

These relations are not presented as predictions but as structural consequences of the RDCC
architecture. They demonstrate that the existence of a finite invariant speed is compatible
with—and may even be illuminated by—the bipartite, CPT-symmetric description of the
Universe. The interpretation developed here remains agnostic about the microscopic origin of
$c$, but it provides a coherent framework in which the emergence of a causal boundary can be
understood as a consequence of global coherence constraints.

A complete assessment of this proposal will require a quantitative treatment of inter-sector
entanglement entropy, its evolution under the dimension-six operator, and the precise mapping
between coherence propagation and sectoral causal structure. The present Companion therefore
serves as a conceptual bridge between the dynamical foundations established in Companions~III–VI
and future work aimed at a more rigorous information-theoretic formulation of RDCC.

\section{Outlook: Physical Consequences}

Interpreting the invariant speed $c$ as an inter-sector coherence boundary leads to several
conceptual and phenomenological consequences that can be explored within the established
RDCC framework without introducing additional dynamical assumptions. These consequences
follow directly from the scaling relations and the decoherence structure developed in this
Companion.

\subsection{Emergent Causal Structure}

The coherence length
\begin{equation}
    \ell_{\mathrm{coh}} = \frac{c}{\alpha_{\mathrm{IR}} H_{\mathrm{IR}}}
\end{equation}
defines the maximal spatial range over which inter-sector correlations remain dynamically
relevant. This induces an effective causal structure that parallels Lieb–Robinson bounds in
quantum many-body systems: the relativistic light cone appears as the saturation limit of
correlation propagation in the underlying bipartite CPT-symmetric state.

In this view, the invariant speed $c$ is not imposed but emerges as the boundary of
coherence-to-correlation transfer. Distances larger than $\ell_{\mathrm{coh}}$ behave as effectively
decohered domains, providing a structural explanation for the universality of causal horizons.

\subsection{Stability of the Cosmological Cycle}

The cycle duration
\begin{equation}
    T_{\mathrm{cycle}} \sim \frac{1}{\alpha_{\mathrm{IR}} H_{\mathrm{IR}}}
\end{equation}
depends only on the frozen infrared relaxation strength and the Hubble rate at freeze-out.
A finite coherence boundary ensures that decoherence proceeds efficiently enough to prevent
the accumulation of long-range inter-sector correlations across cycles, thereby contributing to
the dynamical stability of the RDCC evolution.

This connects the existence of a finite invariant speed to the long-term robustness of the
cyclic background.

\subsection{Connection to Tensor-Sector Phenomenology}

The same dimension-six operator that governs the decay of off-diagonal density-matrix
components also generates the log-periodic modulation of primordial tensor modes. The
coherence-boundary interpretation of $c$ therefore provides a unified structural perspective:
the mechanism that limits correlation propagation also shapes the tensor spectrum.

Future gravitational-wave observations may thus probe the same underlying coherence
structure that determines the emergent value of $c$.

\subsection{Geometric Interpretation}

Although we do not attempt to construct a full emergent metric, the correlation structure
encoded in $\ell_{\mathrm{coh}}$ provides a natural geometric scale. Distances larger than
$\ell_{\mathrm{coh}}$ are effectively decohered, suggesting that aspects of large-scale geometry may
be understood in terms of correlation decay rather than fundamental spacetime primitives.

This perspective is compatible with the global CPT-symmetric structure of RDCC and
requires no modification of Einstein gravity. Instead, it reframes geometric notions as
sectoral projections of a globally coherent quantum state.

\subsection{Future Directions}

A quantitative treatment of correlation propagation in the bipartite state—particularly a
precise definition of inter-sector correlation functions and their evolution under the
dimension-six operator—would allow for a more rigorous derivation of the coherence
boundary and its relation to $c$.

Such developments may clarify whether the emergent causal structure identified here can be
connected to broader approaches in quantum information geometry, entanglement dynamics,
and emergent spacetime. The present Companion therefore serves as a conceptual bridge
between the dynamical foundations established in Companions~III–VI and future work aimed
at a more rigorous information-theoretic formulation of RDCC.


\end{document}
